\section{The diagonal flow}\label{sec:diagonal-flow}

This section is finite-dimensional. It determines the exact asymptotics
of the flow on block-scalar metrics that governs the small spectrum in
Sections~\ref{sec:transfer} and~\ref{sec:complete-spectrum}.

\subsection{Setting and the weight grading}

Let $\Gamma$ be a finite directed acyclic graph with vertex set
$\{1,\ldots,m\}$ and edge set $\Gamma_1$. An edge $\alpha\in\Gamma_1$ from $i$ to
$j$ is written $\alpha=(i\to j)$. Fix masses $M_i>0$ and constants
$c_\alpha>0$, and assume that the underlying undirected graph is connected.
In the geometric application the vertices are the stable factors, an edge
$i\to j$ records a nonzero harmonic entry $\beta_{ij}$ with $F_i$ on the
subbundle side, $M_i=V\rk F_i$, and $c_\alpha=\lVert\beta_{ij}\rVert_{L^2}^2$.

On the hyperplane $\sum_iM_iy_i=0$ of $\mathbb R^m$ put
\begin{equation}\label{eq:diag-flow}
 k_\alpha(y)=c_\alpha e^{2(y_i-y_j)},\qquad
 M_ip_i(y)=\sum_{\alpha\text{ out of }i}k_\alpha
          -\sum_{\alpha\text{ into }i}k_\alpha,\qquad
 \dot y_i=-p_i(y).
\end{equation}
With $d_i=e^{y_i}$ this is the flow $\dot gg^{-1}=-\diag(p_i\id_{F_i})$
for $g=\diag(d_i\id_{F_i})$ studied in Section~\ref{sec:transfer}. Put
$\mathfrak e=\sum_\alpha k_\alpha$, the \emph{total conductance}. The positive eigenvalues
$\nu_1\le\cdots\le\nu_{m-1}$ of the weighted graph form
$\sum_\alpha k_\alpha|x_i-x_j|^2$ on $\{\sum_iM_ix_i=0\}$, with norm
$\sum_iM_i|x_i|^2$, are the \emph{graph eigenvalues} of $y$.

\begin{definition}\label{def:weight-grading}
The \emph{weight grading} is the unique minimizer $w\in\mathbb R^m$ of
$\frac12\sum_iM_iw_i^2$ subject to $w_j-w_i\ge\frac12$ for every edge
$i\to j$. An edge is \emph{tight} if equality holds, and its \emph{slack}
is $\varepsilon_\alpha=2(w_j-w_i)-1\ge0$. The \emph{tight subgraph}
$\mathcal T\subset\Gamma_1$ consists of the tight edges, and its connected components
(on all $m$ vertices, isolated vertices included) are the \emph{tight
components}.
\end{definition}

The constraint set is nonempty, since half the length of a longest path
into each vertex is admissible, and the objective is strictly convex.
The constraints are invariant under adding constants, so
$\sum_iM_iw_i=0$. Since they are linear, there are Lagrange multipliers
$u_\alpha\ge0$, vanishing off $\mathcal T$, with
\begin{equation}\label{eq:kkt}
 M_iw_i=\sum_{\alpha\text{ into }i}u_\alpha
        -\sum_{\alpha\text{ out of }i}u_\alpha .
\end{equation}
Summing over a tight component $\mathcal C$ gives $\sum_{i\in\mathcal C}M_iw_i=0$, because
$u$ is supported on edges inside tight components. A vertex with no tight
edge therefore has $w_i=0$. At least one edge is tight: otherwise every
multiplier would vanish, so $w=0$ by \eqref{eq:kkt}, which is not
admissible.

This is the weight grading of Haiden, Katzarkov, Kontsevich, and Pandit
\cite[Section~3]{HKKP2023}. Their flow \cite[(3.10)]{HKKP2023} is
\eqref{eq:diag-flow} after the substitutions
\[
 x=-2y,\qquad t=2\tau,\qquad v=-2w,\qquad m_i=M_i,
\]
and their tight edges $v_i-v_j=1$ are ours.

\begin{definition}\label{def:nondegenerate}
The pair $(\Gamma,M)$ is \emph{nondegenerate} if some multiplier vector
satisfying \eqref{eq:kkt} is strictly positive on every tight edge.
\end{definition}

In the language of optimization this is strict complementarity. Replacing $M$
by a positive multiple leaves $w$ unchanged and rescales the multipliers,
so for $M_i=V\rk F_i$ nondegeneracy depends only on $\Gamma$ and the ranks.
We then say that $(\Gamma,(r_i))$ is nondegenerate.

This is exactly the hypothesis of \cite[Theorem~3.6]{HKKP2023}, under which
$y=w\log\tau+O(1)$. It fails, for example, for the zigzag
$1\to2\leftarrow3\to4$ with masses satisfying $M_1M_4=M_2M_3$ (compare
\cite[(3.23)--(3.28)]{HKKP2023}),
where iterated logarithms can appear. The equal-rank zigzag of
Theorem~\ref{thm:zigzag-flow} lies on that quadric. In a nondegenerate
graph the tight edges are exactly the edges carrying positive multipliers.

\subsection{Convergence of the remainder}

\begin{theorem}\label{thm:diag-convergence}
Assume $(\Gamma,M)$ is nondegenerate. The tight potential
\[
 \Phi(z)=\sum_iM_iw_iz_i+\frac12\sum_{\alpha\in\mathcal T}c_\alpha e^{2(z_i-z_j)}
\]
has a minimizer $\zeta^*$, unique subject to $\sum_{i\in\mathcal C}M_i\zeta^*_i=0$
for every tight component $\mathcal C$. Every solution of \eqref{eq:diag-flow}
satisfies
\begin{equation}\label{eq:diag-asymptotics}
 y(\tau)=w\log\tau+\zeta^*+\sum_{\mathcal C}Z_{\mathcal C}\mathbf 1_{\mathcal C}+O(\tau^{-\delta})
\end{equation}
for some $\delta>0$ and constants $Z_{\mathcal C}$ with $\sum_{\mathcal C}M_{\mathcal C}Z_{\mathcal C}=0$, where
$M_{\mathcal C}=\sum_{i\in\mathcal C}M_i$ and $\mathbf 1_{\mathcal C}$ is the indicator vector of $\mathcal C$.
The vector $\zeta^*$ depends only on $(\Gamma,M,c)$. The constants $Z_{\mathcal C}$ depend
on the solution.
\end{theorem}

\begin{proof}
Set $s=\log\tau$ for $\tau\ge1$ and $z(s)=y(e^s)-ws$. Then
$\sum_iM_iz_i=0$, and $\tau k_\alpha=c_\alpha e^{2t_\alpha}e^{-\varepsilon_\alpha s}$
with $t_\alpha=z_i-z_j$. Define
\[
 \Psi_s(z)=\sum_iM_iw_iz_i
 +\frac12\sum_\alpha c_\alpha e^{2t_\alpha}e^{-\varepsilon_\alpha s}.
\]
Since $\partial\Psi_s/\partial z_i=M_iw_i+M_i\tau p_i$ and
$dz_i/ds=-\tau p_i-w_i$, the curve $z$ is the gradient flow
$dz/ds=-M^{-1}\nabla\Psi_s(z)$. Because $\partial_s\Psi_s\le0$,
\begin{equation}\label{eq:psi-monotone}
 \frac{d}{ds}\Psi_s(z(s))
 =-\sum_i\frac{(\partial_i\Psi_s)^2}{M_i}
  -\frac12\sum_\alpha\varepsilon_\alpha\tau k_\alpha\le0.
\end{equation}

Fix multipliers $u$ satisfying \eqref{eq:kkt} with $u_\alpha>0$ on $\mathcal T$.
By \eqref{eq:kkt}, $\sum_iM_iw_iz_i=-\sum_\alpha u_\alpha t_\alpha$, so
\[
 \Psi_s(z)=\sum_{\alpha\in\mathcal T}\Bigl(\frac{c_\alpha}2e^{2t_\alpha}
 -u_\alpha t_\alpha\Bigr)
 +\frac12\sum_{\alpha\notin\mathcal T}c_\alpha e^{2t_\alpha}e^{-\varepsilon_\alpha s}.
\]
Each tight term is bounded below by $\frac{u_\alpha}2(1-\log(u_\alpha/c_\alpha))$
and is proper in $t_\alpha$; the other terms are nonnegative. Hence
$\Psi_s(z(s))$ decreases to a finite limit, every tight difference
$t_\alpha$ is bounded along the solution, and \eqref{eq:psi-monotone} gives
\begin{equation}\label{eq:slack-integrable}
 \int_0^\infty\sum_{\alpha\notin\mathcal T}\varepsilon_\alpha\tau k_\alpha\,ds<\infty.
\end{equation}
Within a tight component, $z_i-z_j$ is therefore bounded.

Let $Z_{\mathcal C}(s)=M_{\mathcal C}^{-1}\sum_{i\in\mathcal C}M_iz_i$. In
$\sum_{i\in\mathcal C}\partial_i\Psi_s$ every edge inside $\mathcal C$ cancels, and
$\sum_{i\in\mathcal C}M_iw_i=0$. Edges between different tight components are not
tight, so their slack is at least
$\varepsilon_*=\min_{\alpha\notin\mathcal T}\varepsilon_\alpha>0$. Hence
\[
 \Bigl|\frac{dZ_{\mathcal C}}{ds}\Bigr|
 \le\frac1{M_{\mathcal C}\varepsilon_*}\sum_{\alpha\notin\mathcal T}\varepsilon_\alpha\tau k_\alpha,
\]
which is integrable by \eqref{eq:slack-integrable}. So $Z_{\mathcal C}(s)$ converges,
and $z$ is bounded. Consequently $\tau k_\alpha\le Ce^{-\varepsilon_*s}$ for
$\alpha\notin\mathcal T$. If every edge is tight there is one component, and these
terms are absent.

Let $H$ be the $M$-orthogonal complement of the span of the $\mathbf 1_{\mathcal C}$,
and write $z=\zeta+\sum_{\mathcal C}Z_{\mathcal C}\mathbf 1_{\mathcal C}$ with $\zeta\in H$. The tight
potential satisfies $\Phi(z+\sum_{\mathcal C}\theta_{\mathcal C}\mathbf 1_{\mathcal C})=\Phi(z)$, so
$M^{-1}\nabla\Phi$ takes values in $H$, and
\[
 \frac{d\zeta}{ds}=-M^{-1}\nabla\Phi(\zeta)+R(s),\qquad
 |R(s)|\le Ce^{-\varepsilon_*s}.
\]
On $H$ the map $\zeta\mapsto(t_\alpha)_{\alpha\in\mathcal T}$ is injective, so
$\Phi|_H$ is strictly convex and proper, with a unique minimizer $\zeta^*$,
and its Hessian $2\sum_{\alpha\in\mathcal T}c_\alpha e^{2t_\alpha}(dt_\alpha)^2$ is
uniformly positive on bounded subsets of $H$.

Let $B\subset H$ be a closed ball containing $\zeta^*$ and the bounded curve
$\zeta([0,\infty))$. On $B$, in the norm $|x|_M^2=\sum_iM_ix_i^2$, the
Hessian of $\Phi|_H$ is bounded below by some $c_0>0$. Put
$W=\Phi(\zeta)-\Phi(\zeta^*)\ge0$. For $\zeta\in B$, strong convexity gives
\[
 \tfrac{c_0}2|\zeta-\zeta^*|_M^2\le W\le\tfrac1{2c_0}|\nabla\Phi(\zeta)|_{M^{-1}}^2,
\]
where $|v|_{M^{-1}}^2=\sum_iv_i^2/M_i$. Along the solution,
\[
 \frac{dW}{ds}=-|\nabla\Phi|_{M^{-1}}^2+\langle\nabla\Phi,R\rangle
 \le-\tfrac12|\nabla\Phi|_{M^{-1}}^2+\tfrac12|R|_M^2
 \le-c_0W+C^2e^{-2\varepsilon_*s}.
\]
By Gr\"onwall's inequality, $W(s)\le C'(1+s)e^{-\min(c_0,2\varepsilon_*)s}$,
and hence $|\zeta(s)-\zeta^*|_M=O(e^{-\delta s})$ for every
$\delta<\frac12\min(c_0,2\varepsilon_*)$. Since $z$ is bounded,
$|dZ_{\mathcal C}/ds|\le Ce^{-\varepsilon_*s}$, and so
$|Z_{\mathcal C}(s)-\lim Z_{\mathcal C}|\le Ce^{-\varepsilon_*s}/\varepsilon_*$.
In the variable $\tau=e^s$ this is \eqref{eq:diag-asymptotics}.
\end{proof}

The critical equation of $\Phi$ on $H$ states that
\begin{equation}\label{eq:ustar}
 u^*_\alpha=c_\alpha e^{2(\zeta^*_i-\zeta^*_j)}\qquad(\alpha\in\mathcal T)
\end{equation}
satisfies \eqref{eq:kkt}. It is the unique multiplier vector of this form.
When $\mathcal T$ is a forest, \eqref{eq:kkt} determines $u^*$, which is then
independent of $c$.

\begin{remark}\label{rem:implicit-bias}
After the substitution $\theta=M^{1/2}y$, the flow \eqref{eq:diag-flow} is
gradient flow for the exponential loss $\frac12\sum_\alpha c_\alpha
e^{-\langle x_\alpha,\theta\rangle}$ of the linearly separable vectors
$x_\alpha=2M^{-1/2}(e_j-e_i)$. The weight grading is the corresponding
hard-margin direction, and tight edges are support vectors. For gradient
descent with unit weights, Soudry, Hoffer, Nacson, Gunasekar, and Srebro
prove convergence of the remainder when the support vectors span the data
\cite[Theorems~3 and~4]{Soudry2018}. Here that means there is a single tight
component. Nacson et al.\ \cite[Theorem~10]{Nacson2019} bound the component
of the remainder outside the span of the support vectors without that
assumption. The iterated-logarithm expansion of Soudry et al.\ in the degenerate case
\cite[Theorem~13]{Soudry2018} parallels \cite[Section~5]{HKKP2023}. The
case of several tight components, where the constants $Z_{\mathcal C}$ carry the
dependence on the solution, is the case needed below. The monotone quantity
$\Psi_s$ handles it directly.
\end{remark}

\subsection{Conductances and graph eigenvalues}

Put $a_\alpha=1+\varepsilon_\alpha$ and
$u^0_\alpha=c_\alpha e^{2(\zeta^*_i-\zeta^*_j)}$ for every edge; on tight
edges $u^0_\alpha=u^*_\alpha$. By
Theorem~\ref{thm:diag-convergence},
\begin{equation}\label{eq:edge-asymptotics}
 \tau^{a_\alpha}k_\alpha(y(\tau))\longrightarrow
 k^\infty_\alpha=u^0_\alpha e^{2(Z_{\mathcal C(i)}-Z_{\mathcal C(j)})},
\end{equation}
where $\mathcal C(i)$ is the tight component of $i$. In particular
$\tau k_\alpha\to u^*_\alpha$ on tight edges, whatever the solution.

For an exponent $a$, contract the connected components of the subgraph of
edges with $a_\alpha<a$, adding masses. Join the resulting vertices by the
edges with $a_\alpha=a$ that connect distinct contracted vertices, with
conductances $k^\infty_\alpha$ summed over parallel edges. Let $d_a$ be the number
of positive eigenvalues of this \emph{$a$-quotient graph}. Edges with
$a_\alpha=a$ inside one contracted vertex are omitted; such edges can
occur.

\begin{theorem}\label{thm:diag-spectrum}
Assume $(\Gamma,M)$ is nondegenerate, and let $y$ solve \eqref{eq:diag-flow}.
\begin{enumerate}
\item For each exponent $a$, exactly $d_a$ graph eigenvalues satisfy
 $\nu=\rho\,\tau^{-a}(1+o(1))$, where $\rho$ runs through the positive
 eigenvalues of the $a$-quotient graph, with multiplicity. The
 corresponding spectral subspaces converge to the quotient eigenspaces,
 embedded as vectors constant on the contracted vertices.
\item At $a=1$ the quotient graph is the tight subgraph, which is
 nonempty, with conductances $u^*$ and masses $M_i$. These coefficients are the same for every
 solution. The value $\frac12$ occurs with multiplicity at least the number
 of tight components with at least two vertices. For each such component $\mathcal C$,
 $w\mathbf 1_{\mathcal C}$ is a $\frac12$-eigenvector.
\item For each $a>1$ with $d_a\ge1$, the sum of the coefficients at exponent $a$
 is
 \[
  \sum_{\alpha}k^\infty_\alpha\bigl(M_A^{-1}+M_B^{-1}\bigr),
 \]
 summed over edges with $a_\alpha=a$ that join distinct contracted vertices
 $A\ne B$. As the initial condition varies, this sum is unbounded, and so is
 the largest coefficient at exponent $a$.
\end{enumerate}
\end{theorem}

\begin{proof}
By \eqref{eq:edge-asymptotics}, each conductance lies within a factor
$1+O(\tau^{-\delta})$ of $k^\infty_\alpha\tau^{-a_\alpha}$. A common relative
perturbation of all conductances moves every graph eigenvalue by at most the same
factor, by min--max. The proof of Theorem~\ref{thm:orders} uses only the
ordering of the exponents. With $t=\tau^{-1}$ and real exponents
$a_\alpha$ in place of $2p_{ij}$, it gives the eigenvalue and eigenspace
statements in (1); the $1+o(1)$ factors change each quotient by $o(1)$.

In (2) nothing is contracted at $a=1$, and \eqref{eq:ustar} identifies the
weights. For every tight edge, $w_j-w_i=\frac12$. By \eqref{eq:kkt},
\[
 \sum_{\alpha\in\mathcal T,\;\alpha\ni i}u^*_\alpha(w_i-w_{\alpha(i)})
 =\frac12\Bigl(\sum_{\alpha\text{ into }i}u^*_\alpha
 -\sum_{\alpha\text{ out of }i}u^*_\alpha\Bigr)=\frac12M_iw_i,
\]
where $\alpha(i)$ is the other endpoint. This identity holds row by row, and
the tight Laplacian is block diagonal over tight components. So
$w\mathbf 1_{\mathcal C}$ is a $\frac12$-eigenvector whenever it is nonzero. It is
nonzero when $\mathcal C$ has an edge, and it is $M$-orthogonal to $\mathbf 1_{\mathcal C}$.

For (3), the trace of the Laplacian of the $a$-quotient graph with respect to the
masses is the displayed sum. Its terms are
$u^0_\alpha e^{2(Z_{\mathcal C(i)}-Z_{\mathcal C(j)})}$ with $\mathcal C(i)\ne \mathcal C(j)$, and the linear
forms $Z_{\mathcal C(i)}-Z_{\mathcal C(j)}$ are nonzero on $\{\sum M_{\mathcal C}Z_{\mathcal C}=0\}$. So the sum is
unbounded as $Z$ ranges over this hyperplane.

Every such $Z$ is approximately attained. Start at a late time $s_0$ from
$\zeta^*+\sum_{\mathcal C}\theta_{\mathcal C}\mathbf 1_{\mathcal C}$, with
$\sum_{\mathcal C}M_{\mathcal C}\theta_{\mathcal C}=0$, and put $\Theta=\max|\theta_{\mathcal C}|$.
Let $s_2\in(s_0,\infty]$ be the supremum of the $s$ for which
$|z(\sigma)-z(s_0)|_\infty\le1$ on $[s_0,s]$. On $[s_0,s_2)$, every
difference $t_\alpha$ is at most $2\Theta+2+\max|\zeta^*_i-\zeta^*_j|$. So every
non-tight conductance satisfies $\tau k_\alpha\le C_1e^{4\Theta}e^{-\varepsilon_*s}$,
where $C_1$ does not depend on $(\theta_{\mathcal C})$ or $s_0$. Let $B$ now
be the closed ball in $H$ about $\zeta^*$ containing all $\zeta$ with
$|z-z(s_0)|_\infty\le1$, for $z(s_0)$ of the above form. Its convexity
constant $c_0$ depends only on $\Gamma$, $M$, $c$ and $\Theta$. There:
\begin{itemize}
\item $|dZ_{\mathcal C}/ds|\le C_2e^{4\Theta}e^{-\varepsilon_*s}$, so
 $|Z_{\mathcal C}(s)-\theta_{\mathcal C}|\le C_2e^{4\Theta-\varepsilon_*s_0}/\varepsilon_*$;
\item the transverse part starts at the critical point, $\zeta(s_0)=\zeta^*$.
 The estimate above, with $R=O(e^{4\Theta-\varepsilon_*s})$, gives
 $|\zeta(s)-\zeta^*|_M\le C_3e^{4\Theta-\varepsilon_*s_0}$.
\end{itemize}
The deviation $|z(s)-z(s_0)|_\infty$ is at most
$C_2e^{4\Theta-\varepsilon_*s_0}/\varepsilon_*
+(\min_iM_i)^{-1/2}C_3e^{4\Theta-\varepsilon_*s_0}$. If $s_0$ is so large
that this is below $\frac12$, the deviation stays below $\frac12$ on
$[s_0,s_2)$. By continuity this
forces $s_2=\infty$. Hence
\[
 \lim_{s\to\infty}Z_{\mathcal C}(s)=\theta_{\mathcal C}+O(e^{4\Theta-\varepsilon_*s_0}).
\]
The flow is autonomous, and a shift of time does not change the leading
coefficient of $\tau^{-a}$. So every vector $(\theta_{\mathcal C})$ is, up to an
arbitrarily small error, the vector of limits $Z_{\mathcal C}$ of some
solution.
\end{proof}
